\chapter{Trabajo futuro}\label{ch:decimo-capitulo}
\section{Introducción}
En nuestra investigación hemos dado un paso adelante en la modelización del frenado magnético y el agotamiento del litio para estrellas de tipo solar introduciendo un campo magnético de intensidad variable $B$ y $\amlt$ adaptativo. Nuestro punto de partida fue el hecho de que, a medida que una estrella evoluciona su velocidad de rotación cambia debido a variaciones en su radio y momento angular, al igual que también ocurre con su $\teff$ y $\gsurf$, y por lo tanto, also similar debería de sucederle a $B$ y $\amlt$. Siguiendo esta línea de trabajo, conseguimos reducir el número de parámetros libres asociados a los modelos, además de añadir la capacidad de hacerlos evolucionar. Aún así, los avances aquí mostrados no fueron suficientes para reproducir de manera simultánea los valores actuales para el Sol de la A(Li) y la $\osun$. Por lo tanto, tenemos camino por delante para seguir investigando. 

\section{Línea de trabajo}
Entre las posibles líneas de trabajo para seguir profundizando en las áreas expuestas en esta tesis, y para intentar responder a las cuestiones en ella expuestas y para las que todavía no hemos encontrado respuestas adecuadas, proponemos las siguientes:

\begin{itemize}
    \item En el corto plazo podemos ensanchar la base de simulaciones realizadas hasta este momento, comenzando por extender el rango de masas estelares utilizado en nuestros modelos. Esto nos permitirá ampliar el conjunto de estrellas procedentes de los catálogos Gaia y GES con los que contrastar las medidas arrojadas por nuestras simulaciones. Una mayor muestra con la que comparar nos permitirá investigar hasta qué punto nuestros modelos teóricos pueden ser compatibles con los datos observados, poniendo especial interés en A(Li), $\teff$ y $\ostar$.
    \item Teniendo en cuenta que las estrellas que giran más rápido agotan más litio, se necesitan mecanismos adicionales, como el efecto de bloqueo causado a consecuencia de la interacción con el disco protoestelar, para explicar las observaciones. La consideración de este nuevo proceso en nuestra modelación aportaría un mecanismo adicional de frenado, con lo que acabaríamos teniendo dos actuando a diferentes edades estelares. Sus efectos combinados podrían acabar estableciendo un mecanismo autorregulador que influiría en la evolución del litio y en la intensidad del campo magnético, provocando que se destruyera una menor cantidad de litio y que la intensidad del campo magnético fuera menor, pudiendo así llegar a obtenerse valores en el orden de los que presenta el Sol.
    \item En el corto-medio plazo, y relacionado con el punto anterior, incorporar la interacción con el disco protoestelar parece ser, a priori, una línea de trabajo prometedora. Aún así, hay que hacer constar que el bloqueo con el disco protoestelar sigue siendo un proceso inferido observacionalmente, el cual también depende de parámetros libres. Por lo tanto, aún queda por desarrollar una formulación físicamente realista que dependa únicamente de las propiedades iniciales estelares. Nos encontraremos ante un escenario similar al expuesto para el parámetro $\amlt$, en el que proponemos afrontar la búsqueda de un formalismo que permita reducir el número de parámetros libres a utilizar en favor de leyes de dependencia sobre otros parámetros estelares fundamentes.
    \item En el medio-largo plazo, proponemos abordar las observaciones realizadas sobre estrellas de tipo M, para las que los modelos estándar no son capaces de proporcionar resultados consistentes con las mismas. Acorde al conocimiento que se tiene sobre este tipo de estrellas, se puede realizar una primera clasificación entre aquellas que son totalmente convectivas y las que desarrollan un núcleo radiativo y una zona convectiva por encima de éste. El límite teórico para la transición entre ambas poblaciones se sitúa entorno a 0.3 $\msun$, correspondiente al tipo espectral M4 \citep{Stassun2010}. Por debajo de este umbral la zona convectiva se extendería hasta alcanzar el núcleo de la estrella. Las estructuras estelares correspondientes a los tipos espectrales anteriores y posteriores a M4 se presuponen bastante dispares las unas de las otras y debido a este hecho, se esperan observar diferentes fenómenos que deberían estar conectados, de manera directa o indirecta, con la transición del mencionado umbral de masa. Entre los mecanismos impactados deberían de encontrarse los que gobiernan el período de rotación de la estrella, la intensidad y configuración del campo magnético, así como las escalas temporales evolutivas asociadas a ellos. Los resultados derivados de las simulaciones ayudarían a entender mejor cómo la estructura de una estrella de tipo M, con una determinada masa inicial y composición química, se comporta ante la presencia de un campo magnético durante su evolución.
    \item En paralelo, también sería interesante contactar con el grupo de trabajo detrás del código de evolución estelar MESA. Como parte de esta tesis se ha dotado al mismo de una serie de nuevas características no presentes en la distribución disponible para la comunidad de usuarios. En el desarrollo de estas rutinas se ha invertido una cantidad considerable de recursos, no solo en la forma de tiempo y esfuerzo empleado para su codificación, sino también en el uso de los equipos computacionales de la Universidad durante las fases de pruebas y simulaciones. La posibilidad de poder contribuir a la comunidad open-source aportando parte del código desarrollado para MESA, e incluso de las rutinas de filtrado realizadas en Octave (compatible con Matlab) sobre los catálogs Gaia y GES, sería una forma amable de agradecer el esfuerzo realizado por otros, sin el cual habría sido más difícil realizar nuestra investigación.
\end{itemize}
\section{-- REVISADO HASTA AQUÍ--}
\endinput
